\section{Conclusiones de la fase ejecutada}
La revisión permitió reunir las fuentes y analizar cada inscripción de una persona en un módulo y una presentación. También mostró que una persona puede tener varias inscripciones; esta relación debe conservarse al estudiar la incertidumbre y al separar los datos para evaluar modelos. Se documentaron diferencias entre la fecha de retiro y el resultado final, repeticiones de actividad y registros posteriores al retiro. Los archivos no permiten establecer la causa de todas estas situaciones.

En cada corte se utilizaron datos observados hasta ese día y se excluyeron los retiros que ya habían ocurrido. La comparación mostró que las entregas deben interpretarse junto con el calendario. Por ejemplo, en el día 7 no había evaluaciones cuyo vencimiento estuviera dentro de la ventana, aunque algunos estudiantes ya habían entregado. Por ello, la ausencia de una entrega no siempre indica incumplimiento.

Se encontraron relaciones entre la actividad, los días desde la última interacción y el retiro posterior. Estas relaciones cambian entre cortes, pero también cambian las inscripciones analizadas y el tiempo restante de seguimiento. Para revisar esta situación se repitió parte del análisis con un grupo común. Esa comparación utiliza información posterior y sirve como revisión de los resultados, no como población disponible para una alerta temprana.

El estudio cuenta con fuentes revisadas, indicadores construidos y resultados documentados. Se excluyó el conteo de filas como candidato del modelo y se examinó qué cambia al retirar duplicados, utilizar otra definición del coeficiente de variación (CV) o variar los indicadores del análisis de componentes principales (PCA). Estas comprobaciones ayudan a entender los datos; todavía no demuestran una mejora predictiva. La selección final de variables, la comparación de modelos, SHAP y el prototipo permanecen pendientes. Por tanto, el objetivo general sigue en desarrollo y no se afirma una reducción del abandono ni capacidad comprobada para predecir casos nuevos.

\section{Correspondencia con los objetivos}
La Tabla~\ref{editorial:avance} distingue la evidencia disponible de la terminación de cada objetivo. La comprensión de datos y la ingeniería tienen productos ejecutados; su existencia no implica que se haya fijado el conjunto predictivo final ni evaluado el objetivo general. Los estados se refieren a esta versión, sin asignar porcentajes de avance.
\begingroup\small\setlength{\tabcolsep}{4pt}\renewcommand{\arraystretch}{1.18}
\begin{longtable}{>{\raggedright\arraybackslash}p{.21\linewidth}>{\raggedright\arraybackslash}p{.13\linewidth}>{\raggedright\arraybackslash}p{.59\linewidth}}
\caption{Avance documentado de los objetivos específicos; los nombres se abrevian y su formulación se conserva en el capítulo 2.}\label{editorial:avance}\\
\toprule\textbf{Objetivo} & \textbf{Estado} & \textbf{Evidencia y condición pendiente}\\\midrule\endfirsthead
\toprule\textbf{Objetivo} & \textbf{Estado} & \textbf{Evidencia y condición pendiente}\\\midrule\endhead
\bottomrule\endfoot
1. Identificar y caracterizar variables & Parcial & Capítulo~\ref{cap:caracterizacion}: auditoría y exploración ejecutadas. La selección base de A1.3 se explicita en la Sección~\ref{entrega:seleccionbase}; no depende de entrenar modelos. Quedan por acreditar la procedencia de la descarga y la ratificación del alcance académico.\\
2. Derivar indicadores & Parcial & Capítulo~\ref{cap:ingenieria} y Anexo~\ref{corr:anexodiccionario}: 26 candidatos y conjuntos por corte. A2.5 se limita a entregas y oportunidades; los puntajes permanecen en auditoría. A2.7 cuenta con bases exportadas, pendientes de acordar como entrada al modelado.\\
3. Comparar modelos & Parcial & La etiqueta temporal de A3.1 está implementada, con diferencias documentadas respecto del anteproyecto. El entrenamiento y la evaluación no se han iniciado; el Capítulo~\ref{editorial:modelos} delimita su estado.\\
4. Interpretar y seleccionar el modelo & No iniciado & Capítulo~\ref{editorial:shapestado}: SHAP y selección comparativa pendientes de modelos evaluados. El marco teórico no constituye un resultado de interpretación.\\
5. Diseñar el prototipo & No iniciado & Capítulo~\ref{editorial:prototipo}: pendiente de implementación. Las figuras explicativas de la tesis no equivalen al prototipo de predicciones e interpretabilidad previsto.\\
\end{longtable}\endgroup
Las diferencias de A2.5 y A3.1 requieren la ratificación académica ya señalada. La utilidad educativa se mantiene como propósito de comunicación y apoyo; no se afirma una mejora del aprendizaje o una reducción del abandono mediante intervenciones, cuya evaluación está excluida del alcance.

\section{Continuidad y trabajos futuros}
El siguiente paso es definir cómo se evaluarán los modelos y qué situaciones se espera que puedan predecir: nuevas personas, otras presentaciones o ambas. El protocolo deberá tener en cuenta las inscripciones repetidas y la exploración que ya se hizo de los datos. También deberá fijar cómo comparar ventanas, probabilidades y umbrales. La disponibilidad de las notas y las discrepancias de la fuente mantienen sus limitaciones; se conservan las exclusiones del anteproyecto.